\documentclass[11pt,a4paper]{article}
\usepackage{fontspec,amsmath,amssymb,amsthm,microtype}
\setmainfont{lmroman10-regular.otf}[BoldFont=lmroman10-bold.otf,ItalicFont=lmroman10-italic.otf,BoldItalicFont=lmroman10-bolditalic.otf]
\usepackage[margin=25mm]{geometry}
\usepackage[unicode,hidelinks]{hyperref}
\setlength{\emergencystretch}{3em}
\setlength{\parskip}{0.4em}
\setlength{\parindent}{0pt}
\allowdisplaybreaks[2]
\raggedbottom
\title{\texorpdfstring{Identify recursive sampling laws with iterated Markov transitions}{Identify recursive sampling laws with iterated Markov transitions}}
\author{Omer Deniz Akyildiz\\[0.4em]\normalsize Prepared with AI assistance}
\date{}
\begin{document}\maketitle
Draft mathematical exposition; these notes may contain mistakes.


The path-space law of an iid-noise recursion agrees with the repeated one-step transition \(\mu\mapsto F_\#(\mu\otimes\eta)\). The proof explicitly establishes independence of the current state and the fresh innovation.



Write \(P_h\mu\) for the pushforward of \(\mu\otimes\gamma\) by the update map, where \(\gamma\) is the innovation law. For an independent-noise realization, the current state and fresh innovation have precisely that product law. Hence \(\rho_{k+1}=P_h\rho_k\), with \(\rho_0=\mu\). Prefix measurability and independence of a fresh coordinate justify the product identity, and pushforward probability preservation ensures all recursive laws have mass one.


\section{\texorpdfstring{The state as a function of a finite noise prefix}{The state as a function of a finite noise prefix}}


Evaluate the same recursion using only a finite prefix \((w_0,\ldots,w_{n-1})\). At the successor, recurse on the first \(n\) coordinates and use the final coordinate as the next innovation. This makes finite dependence explicit for the fresh-noise independence proof.



A time-\(n\) state can be evaluated from an input in \(E\times E^{\{0,\ldots,n-1\}}\). Restricting the prefix for the recursive call and reserving its final coordinate for the next update makes this finite dependence explicit. It provides the measurable factorization through past coordinates needed to prove independence from the next innovation.


\section{\texorpdfstring{Agreement with the full recursion}{Agreement with the full recursion}}


Induct on prefix length. Restricting an infinite path to its first \(n\) entries gives the same previous state, and the final finite coordinate is \(w_n\). The successor updates thus agree, identifying finite-prefix and infinite-path constructions pointwise.



At length zero both maps return the initial state. At length \(n+1\), their previous states agree by induction after restricting to the first \(n\) coordinates, and both next innovations are \(w_n\). Applying the common update proves equality. Thus replacing the infinite path by its finite prefix does not alter the algorithm output.


\section{\texorpdfstring{Measurability of the finite-\allowbreak{}prefix state}{Measurability of the finite-prefix state}}


The zero map is the initial projection. At the successor, restriction to the earlier finite coordinates and evaluation of the last coordinate are measurable; combine them with the measurable update and the induction hypothesis. This proves the state depends measurably on a finite prefix.



Restriction of a finite coordinate family and evaluation of its last coordinate are measurable product maps. The induction hypothesis gives measurability of the previous state from the shorter prefix. Pair these two measurable outputs and compose with the update. This proves finite-prefix measurability without any probabilistic assumption and allows independence to be passed through the finite recursion.


\section{\texorpdfstring{Independence of the prefix and the fresh innovation}{Independence of the prefix and the fresh innovation}}


The index sets \(\{0,\ldots,n-1\}\) and \(\{n\}\) are disjoint. Independence of the coordinate family on the iid product source therefore gives independence of the whole prefix vector and the next coordinate. This is stronger than pairwise independence of individual past entries.



On the iid path space, disjoint groups of coordinate evaluations generate independent sigma-algebras. The first group contains all indices less than \(n\), and the second contains only \(n\). Therefore the entire past prefix vector is independent of the fresh innovation. This joint statement is what permits an arbitrary measurable function of the prefix, including the current state, to inherit independence.


\section{\texorpdfstring{The transition acting on a law}{The transition acting on a law}}


Define \(M_\eta(\mu)=F_*(\mu\otimes\eta)\). Its product measure means the innovation is fresh and independent of the current state. The subsequent theorem derives that independence for the canonical recursion, rather than assuming it at each step.



The transition sends \(\mu\) to the law of \(F(X,Z)\) when \(X\sim\mu\), \(Z\sim\eta\), and \(X\perp Z\). Writing it as a product pushforward records these requirements exactly. Knowing only the two marginals would not determine the output if state and innovation were dependent, so the independence verification for the actual recursion is essential.


\section{\texorpdfstring{The iterated Markov laws}{The iterated Markov laws}}


Set \(\mu_0=\mu\), \(\mu_{n+1}=M_\eta(\mu_n)\). These are iterated independent-noise transition laws. The recursion alone does not yet assert that a previously defined sample-path algorithm realizes them.



The law recursion begins with \(\mu\) and repeatedly applies the same independent-noise transition. It defines a candidate sequence of algorithm distributions without constructing sample paths. The next results show that the measurable iid path recursion realizes precisely this sequence. This separates a transition-level definition from the theorem identifying it with an actual noise-driven realization.


\section{\texorpdfstring{The joint law of a state and fresh innovation}{The joint law of a state and fresh innovation}}


The current state is a measurable function of initialization and the first \(n\) noises. Product-source independence and prefix/\allowbreak{}fresh independence make it independent of \(w_n\). Its marginal is the output law of a noise-driven chain, and the innovation marginal is \(\eta\); their joint law is therefore the product of those actual marginals.



The current state factors through initialization and its finite innovation prefix. Initialization is independent of the entire path, and prefix/\allowbreak{}fresh independence separates the next coordinate from that past. Regrouping these facts shows independence of the current state and next innovation. Their joint law is consequently the product of the already defined current output law and \(\eta\).


\section{\texorpdfstring{The successor marginal transition}{The successor marginal transition}}


The pointwise successor is \(F(X_n,w_n)\). Push the proved product joint law of \((X_n,w_n)\) through \(F\); this yields \(\mu_{n+1}=M_\eta(\mu_n)\) for the canonical recursion laws. No later-time state/\allowbreak{}noise independence is assumed.



The successor state is the measurable image of \((X_n,w_n)\) under \(F\). The preceding product-law identity therefore gives its distribution as \(F_\#(\mu_n\otimes\eta)\). This proves the Markov successor relation from the supplied iid source. It does not assume conditional independence as a new property of the recursively constructed state.


\section{\texorpdfstring{Agreement of path-\allowbreak{}space and Markov laws}{Agreement of path-space and Markov laws}}


The canonical path law and the iterated Markov law have the same initial value and the same successor equation. Induction on \(n\) identifies them at every index. This closes the connection between sample-path and transition-law definitions.



Equality at time zero follows from the initial marginal computation. If the two law sequences agree at time \(n\), their common successor transition gives agreement at \(n+1\). This induction identifies the path-space and transition-space formulations for every finite iteration. All later kernel estimates can therefore be transferred to the original finite recursion.


\section{\texorpdfstring{Probability preservation by one transition}{Probability preservation by one transition}}


The independent product of two probability measures has mass one. A measurable pushforward preserves that total mass, so a transition acting on probability measures is again a probability measure.



The product of current-state and innovation probabilities has total mass one. The preimage of the entire output space under a measurable update is the entire product input space. Hence the pushforward transition also has mass one. This conclusion concerns normalization alone and does not assume that the transition preserves second moments or has an invariant distribution.


\section{\texorpdfstring{Probability preservation at every iterate}{Probability preservation at every iterate}}


Induct from the initial probability and apply probability preservation by a measurable Gaussian transition at each successor. This proves normalization of every transition iterate and provides the probability hypothesiss used in later law comparisons.



Assume the initial law is a probability. The one-step normalization theorem shows that each successor law is a probability whenever its predecessor is. Induction proves the property at every finite index. Keeping this fact available is necessary for later expectation inequalities, where the norm of a constant and the comparison between first and second moments use unit total mass.


\section{\texorpdfstring{The noise-\allowbreak{}driven chain has the iterated transition law}{The noise-driven chain has the iterated transition law}}


First identify the independent iid realization with the canonical product-path output law. Then apply the equality of the independent-noise marginal recursion and iterated transition laws. Every realization satisfying the original initialization/\allowbreak{}noise assumptions consequently has the same proved Markov iterate law.



An arbitrary admissible realization first has the same law as the canonical independent path recursion, by equality of its joint input distribution. The canonical path recursion then has the same law as the iterated transition. Combining these two identities yields the claimed Markov-law representation for the supplied process, including all initialization and innovation assumptions used to derive it.


\end{document}
